\chapter{Cypress}

\section{Ce este Cypress?}

Cypress este un framework de testare end-to-end conceput exclusiv pentru aplicații web,
lansat în 2017 de Cypress.io. Spre deosebire de Playwright, care controlează
browserul din exterior printr-un protocol de automatizare, Cypress rulează \textbf{direct
în interiorul browserului}, în același proces JavaScript cu aplicația testată. Această
arhitectură îi conferă acces nativ la DOM, la rețea și la starea internă a aplicației,
fără latența unui strat intermediar.

Cypress suportă \textbf{Chromium} (Chrome, Edge) și \textbf{Firefox}. Suportul pentru
WebKit/Safari este disponibil în versiuni experimentale. Limbajul de scripting este
\textbf{JavaScript sau TypeScript}, iar în cadrul acestui proiect a fost utilizat
\textbf{TypeScript în mod strict}.

\section{Facilități principale}

\subsection{Time Travel Debugging}

Cypress App (interfața grafică interactivă) înregistrează un snapshot al DOM-ului înainte
și după fiecare comandă executată. Trecând cu mouse-ul peste oricare pas din log, interfața
restaurează vizual starea aplicației la acel moment, permițând inspecția elementelor exact
în momentul în care comanda a acționat asupra lor. Aceasta este una din facilitățile cele
mai apreciate de Cypress, eliminând nevoia de a re-rula testul pentru a diagnostica un eșec.

\subsection{Retry-ability automat}

Comenzile Cypress nu returnează imediat o valoare --- ele sunt \textit{lazy} și se re-evaluează
automat până când condiția este îndeplinită sau timeout-ul expiră. De exemplu,
\texttt{cy.get('.element').should('be.visible')} va re-interoga DOM-ul repetat timp de până la
10 secunde (implicit) înainte de a eșua. Această caracteristică elimină necesitatea
așteptărilor explicite în marea majoritate a cazurilor.

\subsection{Interceptarea rețelei cu \texttt{cy.intercept()}}

Cypress oferă un API de interceptare a cererilor HTTP direct la nivelul browserului.
Comanda \texttt{cy.intercept()} permite:
\begin{itemize}
    \item atribuirea unui alias unei cereri și așteptarea ei cu \texttt{cy.wait('@alias')};
    \item inspecția răspunsului (status, body);
    \item modificarea sau blocarea cererilor pentru simularea erorilor.
\end{itemize}

\subsection{Comenzi reutilizabile și helper functions}

TypeScript permite definirea de funcții helper tipizate, utilizabile în orice test.
Proiectul folosește această facilitate extensiv pentru a evita duplicarea codului de
checkout, autentificare și completare formulare.

\subsection{Plugin ecosystem}

Cypress dispune de un ecosistem de plugin-uri pentru extinderea capabilităților native.
Proiectul utilizează \texttt{cypress-real-events}, care simulează evenimente native de
browser (hover, drag, scroll real) --- necesare pentru meniurile dropdown care nu
răspund la evenimentele sintetice Cypress standard.

\subsection{Cypress App --- interfața interactivă}

Comanda \texttt{npm run cy:open} deschide Cypress App, o interfață grafică în care testele
pot fi rulate individual, cu vizualizarea în timp real a fiecărui pas, a snapshot-urilor DOM
și a log-ului de rețea. Aceasta accelerează semnificativ procesul de scriere și depanare a testelor.

\section{Configurarea proiectului}

\subsection{Dependențe}

Dependențele sunt definite în \texttt{package.json}:

\begin{lstlisting}[language=bash, caption={Instalare dependențe Cypress}]
npm install
\end{lstlisting}

Pachetele principale:

\begin{itemize}
    \item \texttt{cypress@15.14.0} --- framework-ul principal;
    \item \texttt{typescript@6.0.3} --- compilatorul TypeScript;
    \item \texttt{cypress-real-events@1.15.0} --- simulare evenimente native de browser.
\end{itemize}

\subsection{Fișierul de configurare}

Configurația relevantă din \texttt{cypress.config.js}:

\begin{lstlisting}[language=JavaScript, caption={cypress.config.js}]
module.exports = defineConfig({
  chromeWebSecurity: false,
  video: false,
  screenshotOnRunFailure: true,
  e2e: {
    baseUrl: 'http://localhost:8080',
    viewportWidth: 1440,
    viewportHeight: 900,
    defaultCommandTimeout: 10000,
    specPattern: 'cypress/e2e/**/*.cy.ts',
  },
});
\end{lstlisting}

\texttt{chromeWebSecurity: false} este necesar pentru a permite accesul cross-origin în
fluxurile de checkout, unde unele resurse sunt încărcate de pe domenii diferite.
\texttt{screenshotOnRunFailure: true} salvează automat un screenshot la fiecare test eșuat,
util pentru diagnosticarea problemelor în rulările headless din CI.

\subsection{Scripturi disponibile}

\begin{itemize}
    \item \texttt{npm test} --- rulare headless completă (\texttt{cypress run});
    \item \texttt{npm run cy:open} --- Cypress App, mod interactiv;
    \item \texttt{npm run cy:run:chrome} --- rulare headed în Chrome;
    \item \texttt{npm run cy:open:local} --- mod interactiv cu \texttt{baseUrl} explicit pe Docker.
\end{itemize}

\section{Implementarea testelor}

\subsection{Structura testelor}

Testele sunt organizate cu \texttt{describe} pentru grupare și \texttt{it} pentru cazuri
individuale, urmând convenția standard Mocha (utilizată intern de Cypress):

\begin{lstlisting}[language=JavaScript, caption={Structura de organizare a testelor}]
describe('NopCommerce Workflows', () => {
  describe('Guest Purchase Journeys', () => {
    it('Guest shopping flow: wishlist -> cart -> checkout success', () => {
      cy.visit('/');
      // ...
    });
  });

  describe('Account-Based Journeys', () => {
    it('Register and place an order in the same session', () => {
      // ...
    });
  });
});
\end{lstlisting}

\subsection{Funcții helper reutilizabile}

O parte semnificativă a codului este extrasă în funcții helper definite la nivelul fișierului,
eliminând duplicarea între teste:

\begin{lstlisting}[language=JavaScript, caption={Helper pentru completarea formularelor de adresă}]
type AddressPayload = {
  firstName: string; lastName: string; email: string;
  city: string; address1: string;
  zipPostalCode: string; phoneNumber: string;
};

const fillAddressForm = (prefix: string, payload: AddressPayload) => {
  cy.intercept('GET', '**/getstatesbycountryid*').as(`getStates_${prefix}`);
  cy.get(`#${prefix}_CountryId`).select('237');
  cy.get(`#${prefix}_StateProvinceId option`)
    .should('have.length.greaterThan', 1);
  cy.get(`#${prefix}_StateProvinceId`).select(1);
  cy.get(`#${prefix}_FirstName`).clear().type(payload.firstName);
  // ...
};
\end{lstlisting}

Funcția \texttt{fillAddressForm} este apelată cu un prefix diferit pentru fiecare tip de
adresă (\texttt{BillingNewAddress}, \texttt{ShippingNewAddress}, \texttt{Address}),
construind dinamic ID-urile câmpurilor.

\subsection{Interceptarea cererilor AJAX}

Interceptarea cererilor AJAX în Cypress se face prin \texttt{cy.intercept()} combinat cu
\texttt{cy.wait()}. Exemplul următor verifică că răspunsul la filtrarea produselor
returnează HTTP 200:

\begin{lstlisting}[language=JavaScript, caption={Interceptare AJAX la aplicarea filtrului}]
cy.intercept('**/*').as('ajaxFilter');
cy.get("input[id='attribute-option-9']").check({ force: true });
cy.wait('@ajaxFilter').its('response.statusCode').should('eq', 200);
\end{lstlisting}

Pentru cereri mai specifice, se folosește o expresie regulată:

\begin{lstlisting}[language=JavaScript, caption={Interceptare cerere de actualizare atribute produs}]
cy.intercept(/productdetails_attributechange/i).as('attrChange1');
cy.get("input[name='product_attribute_3'][value='7']").check({ force: true });
cy.wait('@attrChange1');
\end{lstlisting}

\subsection{Verificarea prețurilor dinamice}

Funcția utilitară \texttt{toNumber()} curăță textul afișat și îl convertește la număr,
permițând comparații matematice:

\begin{lstlisting}[language=JavaScript, caption={Extragere și verificare preț dinamic}]
const toNumber = (text: string) =>
  Number(text.replace(/[^\d.,-]/g, '').replace(/,/g, ''));

cy.get('div.product-price span').invoke('text').then((basePriceText) => {
  const basePrice = toNumber(basePriceText);

  // selectie componente + asteptare AJAX
  cy.intercept(/productdetails_attributechange/i).as('attrChange');
  cy.get("input[name='product_attribute_3'][value='7']").check({ force: true });
  cy.wait('@attrChange');

  cy.get('div.product-price span').invoke('text').then((newPriceText) => {
    const newPrice = toNumber(newPriceText);
    expect(newPrice).to.be.greaterThan(basePrice);
  });
});
\end{lstlisting}

\subsection{Soft assertions}

Cypress nu oferă soft assertions nativ. Acestea se implementează prin \texttt{cy.get('body').then()}
care permite inspecția condiționată a DOM-ului fără a opri execuția testului:

\begin{lstlisting}[language=JavaScript, caption={Soft assert pentru cuponul invalid}]
cy.get('body').then(($body) => {
  if ($body.find('div.message-failure').length > 0) {
    cy.get('div.message-failure').first().invoke('text').then((failureText) => {
      const normalized = failureText.trim().toLowerCase();
      if (normalized.length > 0) {
        expect(normalized).to.contain('coupon');
      } else {
        cy.log('Coupon container present but empty, skipping assertion.');
      }
    });
  } else {
    cy.log('Coupon error banner did not appear, continuing as soft assertion.');
  }
});
\end{lstlisting}

\subsection{Hover cu cypress-real-events}

Meniurile dropdown ale NopCommerce se activează la hover. Evenimentele sintetice Cypress
(\texttt{cy.trigger('mouseover')}) nu declanșează CSS \texttt{:hover} în toate cazurile.
Plugin-ul \texttt{cypress-real-events} rezolvă aceasta prin simularea unui eveniment nativ:

\begin{lstlisting}[language=JavaScript, caption={Hover nativ pentru meniul dropdown}]
cy.get("a.menu__link[href='/computers']").realHover();
cy.get("a[href='/notebooks']").click({ force: true });
\end{lstlisting}

\subsection{Gestionarea metodelor de plată variabile}

Pe instanța demo, metodele de plată disponibile pot varia. Funcția helper
\texttt{selectPaymentMethodAndContinue()} gestionează acest caz prin inspecția DOM-ului
înainte de selecție:

\begin{lstlisting}[language=JavaScript, caption={Selecție robustă a metodei de plată}]
const selectPaymentMethodAndContinue = () => {
  cy.get('body').then(($body) => {
    const creditCardMethod = $body.find(
      "input[type='radio'][name='paymentmethod'][value*='Manual']"
    );
    const allMethods = $body.find(
      "input[type='radio'][name='paymentmethod']"
    );

    if (creditCardMethod.length > 0) {
      cy.wrap(creditCardMethod.first()).check({ force: true });
      return;
    }
    if (allMethods.length > 0) {
      cy.wrap(allMethods.first()).check({ force: true });
    }
  });

  cy.get('.payment-method-next-step-button:visible').click();
};
\end{lstlisting}

\section{Avantaje}

\begin{itemize}
    \item \textbf{Developer experience superior}: Time Travel Debugging, hot reload în Cypress App
          și feedback vizual în timp real accelerează dramatic procesul de scriere a testelor.
    \item \textbf{Retry-ability built-in}: eliminarea așteptărilor explicite face testele mai
          concise și mai reziliente la variații de latență.
    \item \textbf{\texttt{cy.intercept()} intuitiv}: API-ul de interceptare este simplu și expresiv,
          cu suport pentru aliasuri și aserțiuni înlănțuite.
    \item \textbf{TypeScript nativ}: tipizarea strictă detectează erori de selector sau de tip
          la momentul compilării, nu la runtime.
    \item \textbf{Setup rapid}: \texttt{npm install} + \texttt{npx cypress open} este suficient
          pentru a începe --- fără instalare separată de binarele de browser.
    \item \textbf{Screenshots automate la eșec}: documentarea eșecurilor fără configurare
          suplimentară.
\end{itemize}

\section{Dezavantaje}

\begin{itemize}
    \item \textbf{Suport browser limitat}: Cypress suportă nativ doar Chromium și Firefox.
          Suportul WebKit/Safari este experimental și nu acoperă toate cazurile de utilizare.
    \item \textbf{JavaScript/TypeScript exclusiv}: nu există suport pentru Python, Java sau C\#,
          limitând adoptarea în echipe cu alt stack tehnic.
    \item \textbf{Fără stealth mode}: Cypress nu are un echivalent al \texttt{playwright-stealth}.
          Testele pe instanțe protejate de Cloudflare necesită un mediu local (Docker), ceea
          ce crește complexitatea infrastructurii.
    \item \textbf{Arhitectură single-tab}: Cypress rulează într-o singură filă de browser;
          scenariile care implică mai multe tab-uri sau ferestre sunt dificil de testat.
    \item \textbf{Soft assertions manuale}: similar cu Playwright, absența unui mecanism nativ
          obligă la construcții \texttt{cy.get('body').then()} verbale.
    \item \textbf{Asincronicitate implicită}: comenzile Cypress sunt înlănțuite asincron
          (\textit{command queue}), iar accesul la valori în afara callback-urilor
          \texttt{.then()} este o sursă frecventă de confuzie pentru utilizatorii noi.
\end{itemize}
